%============================================================
%  Iterated Symmetric Centipede: Two-Population Analysis
%  Replicates structure of Lazaridis & Kehagias (2026) with
%  two populations instead of one.
%============================================================
\documentclass[12pt,a4paper]{article}

\usepackage[utf8]{inputenc}
\usepackage[T1]{fontenc}
\usepackage{amsmath,amssymb,amsthm}
\usepackage{graphicx}
\usepackage{booktabs}
\usepackage[margin=2.5cm]{geometry}
\usepackage{caption}
\usepackage{subcaption}
\usepackage{float}
\usepackage{hyperref}
\usepackage{array}
\usepackage{multirow}

\newtheorem{definition}{Definition}[section]
\newtheorem{proposition}{Proposition}[section]
\newtheorem{corollary}{Corollary}[section]
\newtheorem{example}{Example}[section]
\newtheorem{remark}{Remark}[section]

\graphicspath{{figures/}}

\title{\textbf{Iterated Symmetric Centipede:\\
Two-Population Learning and Evolutionary Games}}
\author{}
\date{2026}

\begin{document}

\maketitle

%------------------------------------------------------------
\begin{abstract}
We study the Iterated Symmetric Centipede (ISC) game under a
\emph{two-population} framework.  The ISC involves repeated plays of the
One-shot Symmetric Centipede (OSC), in which the first-moving player is chosen
at random.  We consider two separate populations, each using the strategies
All-$M$, All-$1$, and Grim, with every player from Population~1 playing against
every player from Population~2 (no within-population interaction).  We first
identify the Nash equilibria of the two-population ISC bimatrix game
$(C, C^\top)$, which includes both symmetric and \emph{asymmetric} equilibria
not reachable by one-population dynamics.  We then study large-population
evolution via the Two-Population Replicator Dynamics~(B.1)--(B.2) and
small-population evolution via a two-population Markov chain using Pairwise
Proportional Imitation~(PPI).  We show that the threshold $p^*=2/3$ governing
the stability of the all-defector equilibrium is preserved from the
one-population model, and that the two-population framework generically
produces asymmetric long-run outcomes in which one population cooperates fully
while the other plays Grim.
\end{abstract}
%------------------------------------------------------------

\tableofcontents
\newpage

%============================================================
\section{Introduction}
%============================================================

This paper extends the framework of Lazaridis and Kehagias~\cite{LK2026} by
replacing the single evolving population with \emph{two} populations that play
exclusively against each other.  The motivation is twofold.

First, one-population replicator dynamics can only converge to \emph{symmetric}
Nash equilibria~\cite{Sandholm2010}. The ISC possesses non-symmetric NE (e.g.\
Population~1 playing All-$M$ while Population~2 plays Grim), which are
strategically meaningful but inaccessible to one-population dynamics.  The
two-population framework recovers these.

Second, many real interactions are naturally inter-group rather than
intra-group: buyers vs.\ sellers, incumbents vs.\ entrants, immune systems vs.\
pathogens.  A two-population model captures this structure directly.

The paper is structured to match~\cite{LK2026} section by section.
Section~\ref{sec:prelim} reviews the classic Centipede and OSC.
Section~\ref{sec:ISC} defines the two-population ISC.
Section~\ref{sec:C} derives the payoff matrix.
Section~\ref{sec:NE} identifies all Nash equilibria, including the new
asymmetric ones.
Section~\ref{sec:rep} presents the two-population replicator dynamics and
stability analysis.
Section~\ref{sec:ESS} discusses evolutionary stability.
Section~\ref{sec:markov} analyses the two-population Markov chain.
Section~\ref{sec:learning} connects to learning theory.
Section~\ref{sec:conc} concludes.

%============================================================
\section{Preliminaries}
\label{sec:prelim}
%============================================================

\subsection{Classic Centipede}

The Centipede game~\cite{Rosenthal1981} is an extensive-form game of perfect
information on a path of $M$ non-terminal vertices.
Player~1 moves at odd turns and Player~2 at even turns.  The total payoff at
turn $t$ is $c_t = t$.  When the active player terminates at turn~$t$, the
payoffs are $(p\,c_t,\,(1-p)\,c_t)$ with terminator's share
$p\in(\tfrac{1}{2},1]$.  We define the strategy family
\[
  \rho_m = \text{``terminate at the first available turn at or after stage~}m\text{''.}
\]
Backward induction yields the unique subgame-perfect equilibrium $(\rho_1,\rho_1)$,
producing payoffs $(p,1-p)$.  Both players would fare better under $(\rho_M,\rho_M)$,
which produces $(pM,(1-p)M)$.

\subsection{One-Shot Symmetric Centipede (OSC)}

To remove asymmetry we randomise the first mover: each player moves first with
probability $\tfrac{1}{2}$.  The resulting normal-form game has
$M\times M$ payoff matrix $A$ for the row player:
\begin{equation}\label{eq:A}
  A(i,j)=\begin{cases}
    i & \text{if }i=j,\\
    (2i+1)\,p & \text{if }i<j,\\
    (2j+1)\,(1-p) & \text{if }i>j.
  \end{cases}
\end{equation}
Four values used throughout:
$A_{MM}=M$, $A_{11}=1$, $A_{M1}=3(1-p)$, $A_{1M}=3p$.

%============================================================
\section{Two-Population ISC}
\label{sec:ISC}
%============================================================

\subsection{Setup}

The ISC consists of $T$ repeated OSC plays with full history observation.
In the two-population model:
\begin{itemize}
  \item \textbf{Population 1} has $N_1$ players, each using one strategy from
        $\{\sigma_1,\sigma_2,\sigma_3\}$.
  \item \textbf{Population 2} has $N_2$ players, each using one strategy from
        $\{\sigma_1,\sigma_2,\sigma_3\}$.
  \item Every Population-1 player plays \emph{exactly once} against every
        Population-2 player.  \textbf{No within-population games occur.}
\end{itemize}

\subsection{Strategies}

\begin{definition}[ISC Strategies]\label{def:strats}
\hfill
\begin{enumerate}
  \item $\sigma_1$ (\textbf{All-$M$}): play $\rho_M$ in every round.
  \item $\sigma_2$ (\textbf{All-$1$}): play $\rho_1$ in every round.
  \item $\sigma_3$ (\textbf{Grim}): play $\rho_M$ until the opponent plays
        $\rho_m$ with $m<M$; then switch permanently to $\rho_1$.
\end{enumerate}
\end{definition}

\subsection{Bimatrix Game}

Since all strategies are available to both populations and payoffs depend only
on (own strategy, opponent's strategy), the game is symmetric:
Population~1's payoff matrix is $C$ (defined below) and Population~2's payoff
matrix is $C^\top$.  We therefore write the bimatrix game as $(C,C^\top)$.

%============================================================
\section{The ISC Payoff Matrix}
\label{sec:C}
%============================================================

\subsection{Derivation}

$C(i,j)$ is the payoff to a strategy-$i$ player in Population~1 facing a
strategy-$j$ player from Population~2 over $T$ OSC rounds.

\begin{proposition}[ISC Payoff Matrix]\label{prop:C}
\begin{equation}\label{eq:C}
C = \begin{pmatrix}
T M & 3(1-p)\,T & T M \\[4pt]
3p\,T & T & 3p+(T-1) \\[4pt]
T M & 3(1-p)+(T-1) & T M
\end{pmatrix}.
\end{equation}
\end{proposition}

\begin{proof}
All-$M$ never plays a sub-$M$ action, so it never triggers Grim and Grim never
triggers against it.  Entries $C(1,\cdot)$, $C(\cdot,1)$, and $C(3,3)$ all
equal $TM$.  For $C(2,3)$: in round~1 All-$1$ faces Grim's $\rho_M$, earning
$A_{1M}=3p$; Grim triggers; rounds $2,\ldots,T$ are mutual $\rho_1$, earning
$A_{11}=1$ each.  Similarly $C(3,2)=A_{M1}+(T-1)=3(1-p)+(T-1)$.  Remaining
entries follow directly from Definition~\ref{def:strats}.
\end{proof}

\subsection{Numerical Values ($M=4$, $T=10$)}

\begin{table}[H]
\centering
\caption{Payoff matrix $C$ for $p=3/4$, $M=4$, $T=10$.
         Population~2 payoff matrix is $C^\top$.}
\label{tab:C_p075}
\renewcommand{\arraystretch}{1.2}
\begin{tabular}{lccc}
\toprule
           & vs All-$M$ & vs All-$1$ & vs Grim \\
\midrule
All-$M$    & 40 & 7.5  & 40    \\
All-$1$    & 22.5 & 10 & 11.25 \\
Grim       & 40 & 9.75 & 40    \\
\bottomrule
\end{tabular}
\end{table}

\begin{table}[H]
\centering
\caption{Payoff matrix $C$ for $p=3/5$, $M=4$, $T=10$.}
\label{tab:C_p060}
\renewcommand{\arraystretch}{1.2}
\begin{tabular}{lccc}
\toprule
           & vs All-$M$ & vs All-$1$ & vs Grim \\
\midrule
All-$M$    & 40 & 12  & 40    \\
All-$1$    & 18 & 10  & 10.8  \\
Grim       & 40 & 10.2 & 40   \\
\bottomrule
\end{tabular}
\end{table}

\begin{remark}[Symmetry of $C$]
$C(1,3)=C(3,1)=C(1,1)=C(3,3)=TM$ and $C(3,2)=3(1-p)+(T-1)\approx T$,
$C(2,3)=3p+(T-1)\approx T$ for large~$T$.  The game is symmetric:
$C^\top(j,i)=C(j,i)$, confirming the bimatrix representation $(C,C^\top)$.
\end{remark}

%============================================================
\section{Nash Equilibria}
\label{sec:NE}
%============================================================

\begin{definition}[Two-Population Nash Equilibrium]\label{def:NE}
A pair of mixed strategies $(\hat x, \hat y)$ with $\hat x,\hat y\in\Delta_3$
is a \emph{Nash equilibrium} of the bimatrix game $(C,C^\top)$ iff
\[
  \forall x\in\Delta_3:\; \hat x^\top C\hat y \geq x^\top C\hat y,\qquad
  \forall y\in\Delta_3:\; \hat y^\top C^\top\hat x \geq y^\top C^\top\hat x.
\]
\end{definition}

\noindent Note that $\hat y^\top C^\top\hat x = \hat x^\top C\hat y$, so both
players share the same equilibrium payoff.

\begin{proposition}[Asymmetric NE Family]\label{prop:NEfamily}
For every $M\in\{3,4,\ldots\}$ and large enough $T$, the following are NE of
the two-population ISC for any $a',a''\in[0,1]$:
\begin{equation}\label{eq:NEfamily}
  \bigl((1-a',\,0,\,a'),\;(1-a'',\,0,\,a'')\bigr).
\end{equation}
\end{proposition}

\begin{proof}
Let $\hat x=(1-a',0,a')$ and $\hat y=(1-a'',0,a'')$.  Since $C(i,j)=TM$ for all
$i,j\in\{1,3\}$, every strategy in the support of $\hat x$ or $\hat y$ earns
$TM$ against the opponent mixture.  The mean payoff is $TM$.  A deviation to
$\sigma_2$ earns at most $\max(C(2,1),C(2,3))=\max(3pT,\,3p+(T-1))<TM$ for $M\geq 2$.
Hence no profitable deviation exists for either population, regardless of the
values of $a'$ and $a''$.
\end{proof}

\begin{remark}[New equilibria vs.\ one-population case]
Proposition~\ref{prop:NEfamily} includes \emph{asymmetric} NE such as
$(e_1,e_3)$ (Pop~1 all All-$M$, Pop~2 all Grim) and $(e_3,e_1)$.  These cannot
be expressed as symmetric NE and are therefore inaccessible to
one-population dynamics (which can only locate symmetric equilibria).
\end{remark}

\begin{proposition}[Characterisation of Symmetric NE]\label{prop:symNE}
For every $M\in\{3,4,\ldots\}$ and large enough $T$, the following are the only
symmetric NE $(\hat x,\hat x)$ of the two-population ISC:
\begin{enumerate}
  \item $((1-a',0,a'),(1-a',0,a'))$ for all $a'\in[0,1]$.
  \item $((0,1,0),(0,1,0))$ if and only if $p\geq\tfrac{2}{3}$.
  \item When $p\geq\tfrac{2}{3}$, the mixed NE
        \[\left(\!\left(\frac{3p-2}{M-2},\frac{M-3p}{M-2},0\right),
                \left(\frac{3p-2}{M-2},\frac{M-3p}{M-2},0\right)\!\right)\]
        and
        \[\left(\!\left(0,\frac{T(M-1)+1-3p}{T(M-1)-1},\frac{3p-2}{T(M-1)-1}\right),
                \left(0,\frac{T(M-1)+1-3p}{T(M-1)-1},\frac{3p-2}{T(M-1)-1}\right)\!\right).\]
\end{enumerate}
These are exactly the symmetric NE from~\cite{LK2026}, Proposition~4.3.
\end{proposition}

\begin{proof}
For a symmetric NE $(\hat x,\hat x)$, both populations face the same mixture and
the problem reduces to the one-population game.  The symmetric NE of the
bimatrix game $(C,C^\top)$ coincide with the NE of the one-population game with
matrix $C$.  The result then follows from~\cite{LK2026}, Proposition~4.3.
\end{proof}

\begin{example}[$M=4$, $T=10$, $p=3/4$]
The NE include: all pairs from~\eqref{eq:NEfamily} (e.g.\ $(e_1,e_3)$,
$(e_3,e_1)$, $(e_1,e_1)$, $(e_3,e_3)$, and every mixed combination),
plus the symmetric mixed NE
$\bigl((\tfrac{1}{8},\tfrac{7}{8},0),(\tfrac{1}{8},\tfrac{7}{8},0)\bigr)$
and
$\bigl((0,\tfrac{115}{116},\tfrac{1}{116}),(0,\tfrac{115}{116},\tfrac{1}{116})\bigr)$
from Proposition~\ref{prop:symNE}.
\end{example}

%============================================================
\section{Two-Population Replicator Dynamics}
\label{sec:rep}
%============================================================

\subsection{Model}

Let $x(t)=(x_1,x_2,x_3)$ and $y(t)=(y_1,y_2,y_3)$ denote the strategy
frequency vectors of Population~1 and Population~2 at time $t$, respectively.
Because there are no within-population games, the fitness of strategy $m$ in
Population~1 is $(Cy)_m$ and in Population~2 it is $(x^\top C)_m=(C^\top x)_m$.
Both populations share the \emph{same} mean payoff scalar $x^\top Cy$.

The two-population replicator equations are:
\begin{align}
  \forall m:\quad \frac{dx_m}{dt} &= x_m\bigl((Cy)_m - x^\top Cy\bigr),
  \label{eq:rep1}\\[4pt]
  \forall m:\quad \frac{dy_m}{dt} &= y_m\bigl((x^\top C)_m - x^\top Cy\bigr).
  \label{eq:rep2}
\end{align}
The system~\eqref{eq:rep1}--\eqref{eq:rep2} conserves the simplices
$\Delta_3^x=\{x\geq 0,\sum x_m=1\}$ and $\Delta_3^y=\{y\geq 0,\sum y_m=1\}$
separately.

\begin{remark}[Connection to one-population dynamics]
Setting $y=x$~and requiring consistency yields $dx_m/dt=x_m[(Cx)_m-x^\top Cx]$,
the standard one-population replicator.  Hence the one-population dynamics is
the restriction of~\eqref{eq:rep1}--\eqref{eq:rep2} to the diagonal $x=y$.
\end{remark}

\subsection{Rest Points}

\begin{proposition}[Rest Points]\label{prop:rest2pop}
For every $M\in\{3,4,\ldots\}$ and every $T$, the following are admissible rest
points of~\eqref{eq:rep1}--\eqref{eq:rep2}:
\begin{enumerate}
  \item For all $a,b\in[0,1]$:
        \begin{equation}\label{eq:coop_face}
          \bigl((1-a,\,0,\,a),\;(1-b,\,0,\,b)\bigr).
        \end{equation}
  \item $\bigl((0,1,0),\,(0,1,0)\bigr)$ for all $p$.
  \item When $p\geq\tfrac{2}{3}$, the symmetric mixed rest points from
        Proposition~\ref{prop:symNE}.
\end{enumerate}
The set~\eqref{eq:coop_face} is a two-dimensional non-isolated equilibrium set.
\end{proposition}

\begin{proof}
For~\eqref{eq:coop_face}: if $x_2=y_2=0$ then $(Cy)_m=TM$ for
$m\in\{1,3\}$ (since $C(1,j)=C(3,j)=TM$ for $j\in\{1,3\}$).  Hence
$(Cy)_m=x^\top Cy=TM$ and $\dot x_m=0$.  Similarly $(x^\top C)_m=TM$ for
$m\in\{1,3\}$ and $\dot y_m=0$.  The set is two-dimensional (parametrised
by $a$ and $b$ independently).

For $((0,1,0),(0,1,0))$: $(Cy)_2=C(2,2)=T=x^\top Cy$ and
$(x^\top C)_2=C(2,2)=T$.  All derivatives are zero.

Symmetric mixed rest points follow by the same argument as
Proposition~4.1 in~\cite{LK2026}.
\end{proof}

\subsection{Stability of Rest Points}

\begin{proposition}[Stability of the Cooperative Face]\label{prop:stabcoop}
For every $M\geq 4$ and large enough $T$, the entire face
$\mathcal{F}=\{(x,y):x_2=y_2=0\}$ is a stable non-isolated equilibrium set.
Specifically, for any $(x^*,y^*)\in\mathcal{F}$ and perturbations
$\varepsilon_x>0$ (All-$1$ entering Pop~1) and $\varepsilon_y>0$
(All-$1$ entering Pop~2), the Jacobian eigenvalues in the $x_2$ and $y_2$
directions are both strictly negative.
\end{proposition}

\begin{proof}
The eigenvalue governing All-$1$ invasion of Population~1 is
\[
  \lambda_{x_2} = (Cy^*)_2 - x^{*\top}Cy^*
  = \bigl[b\,C(2,1)+(1-b)\,C(2,3)\bigr] - TM,
\]
where $y^*=(1-b,0,b)$.  For large~$T$:
$(Cy^*)_2\approx b\cdot 3pT+(1-b)\cdot T=T[1+(3p-1)b]$.
Hence
\[
  \lambda_{x_2}\approx T[1+(3p-1)b]-TM = T[1-M+(3p-1)b].
\]
Since $M\geq 4$ and $(3p-1)b\leq 2$ for $b\leq 1$, $p\leq 1$:
$1-M+(3p-1)b \leq 1-4+2=-1<0$.  By symmetry $\lambda_{y_2}<0$ for the same
reason.  Hence $\mathcal{F}$ is stable in both transverse directions.
\end{proof}

\begin{remark}
Proposition~\ref{prop:stabcoop} shows the cooperative face $\mathcal{F}$ is
\emph{always} stable for $M\geq 4$, regardless of~$p$.  This contrasts with
the one-population model where the corresponding face was only neutrally stable
and not asymptotically stable.  The key difference is that in the two-population
model the All-$1$ mutant in Population~1 faces only Population~2's cooperation,
giving it payoff $\approx T < TM$, while in the one-population model the
mutant also encounters the other one-population players at a rate that can
sometimes be favourable.
\end{remark}

\begin{proposition}[Stability of All-$1$ Equilibrium]\label{prop:stabA1}
For every $M\in\{3,4,\ldots\}$ and large enough $T$:
\begin{enumerate}
  \item If $p>\tfrac{2}{3}$: $((0,1,0),(0,1,0))$ is asymptotically stable.
  \item If $p\leq\tfrac{2}{3}$: $((0,1,0),(0,1,0))$ is unstable.
\end{enumerate}
\end{proposition}

\begin{proof}
The eigenvalues of the Jacobian of~\eqref{eq:rep1}--\eqref{eq:rep2} at
$((0,1,0),(0,1,0))$ in the directions of each strategy $m\neq 2$ are:
\begin{align*}
  \lambda_{x_1} &= C(1,2)-C(2,2) = 3(1-p)T-T = T(2-3p),\\
  \lambda_{x_3} &= C(3,2)-C(2,2) = 3(1-p)+(T-1)-T = 2-3p,\\
  \lambda_{y_1} &= (C^\top e_2)_1 - e_2^\top Ce_2 = C(1,2)-T = T(2-3p),\\
  \lambda_{y_3} &= C(3,2)-T = 2-3p.
\end{align*}
All four eigenvalues have sign $\text{sgn}(2-3p)$.  For $p>2/3$: negative
$\Rightarrow$ asymptotically stable.  For $p<2/3$: positive $\Rightarrow$
unstable.  For $p=2/3$: a higher-order perturbation analysis (analogous to
Proposition~5.2 in~\cite{LK2026}) shows instability.
\end{proof}

\begin{proposition}[Stability of Asymmetric Pure Equilibria]\label{prop:stabAsym}
For every $M\geq 4$ and large enough $T$, the asymmetric pure equilibria
$(e_1,e_3)$ and $(e_3,e_1)$ are \emph{neutrally stable}: the Jacobian
eigenvalues are all non-positive, with zero eigenvalues in the directions
$(e_1-e_3)^x$ and $(e_1-e_3)^y$ corresponding to the neutrality within the
cooperative face.
\end{proposition}

\begin{proof}
At $(e_1,e_3)$: $(Ce_3)_m=C(m,3)$, which equals $TM$ for $m\in\{1,3\}$
and $C(2,3)\approx T<TM$ for $m=2$.  Eigenvalues:
$\lambda_{x_1}=C(1,3)-TM=0$, $\lambda_{x_2}=C(2,3)-TM<0$,
$\lambda_{x_3}=C(3,3)-TM=0$.
$(e_1^\top C)_n=C(1,n)$: $\lambda_{y_1}=C(1,1)-TM=0$,
$\lambda_{y_2}=C(1,2)-TM<0$, $\lambda_{y_3}=C(1,3)-TM=0$.
All eigenvalues are $\leq 0$; the zeros reflect neutral drift within
$\mathcal{F}$.
\end{proof}

\subsection{Phase Portraits}

We visualise the dynamics on two side-by-side ternary simplices: the left panel
shows Population~1's trajectory and the right panel shows Population~2's
trajectory for the same integrated orbit.  Both populations start from the same
initial condition and are integrated jointly under~\eqref{eq:rep1}--\eqref{eq:rep2}.

\begin{figure}[H]
  \centering
  \begin{subfigure}[b]{0.48\textwidth}
    \includegraphics[width=\textwidth]{rep_phase_p075}
    \caption{$p=3/4$: both populations converge to the cooperative face or
             the All-$1$ vertex.}
    \label{fig:phase_p075}
  \end{subfigure}
  \hfill
  \begin{subfigure}[b]{0.48\textwidth}
    \includegraphics[width=\textwidth]{rep_phase_p060}
    \caption{$p=3/5$: All-$1$ is unstable; both converge to the cooperative
             face, possibly at different points (asymmetric equilibrium).}
    \label{fig:phase_p060}
  \end{subfigure}
  \caption{Phase portraits of~\eqref{eq:rep1}--\eqref{eq:rep2}, $M=4$, $T=10$.
           \emph{Left}: Population~1 simplex (vertices: All-$M$ bottom-left,
           All-$1$ bottom-right, Grim top).  \emph{Right}: Population~2 simplex
           (same layout).  Circle = initial condition; square = final state.}
  \label{fig:phase_both}
\end{figure}

\begin{figure}[H]
  \centering
  \begin{subfigure}[b]{0.48\textwidth}
    \includegraphics[width=\textwidth]{rep_timeseries_p075}
    \caption{$p=3/4$.}
  \end{subfigure}
  \hfill
  \begin{subfigure}[b]{0.48\textwidth}
    \includegraphics[width=\textwidth]{rep_timeseries_p060}
    \caption{$p=3/5$.}
  \end{subfigure}
  \caption{Time series of~\eqref{eq:rep1}--\eqref{eq:rep2}, $M=4$, $T=10$.
           Solid lines: Population~1.  Dashed lines: Population~2.
           Colours: blue = All-$M$, red = All-$1$, green = Grim.
           Initial condition: $x_0=(0.6,0.2,0.2)$, $y_0=(0.1,0.1,0.8)$.}
  \label{fig:rep_timeseries}
\end{figure}

\subsection{Parameter Sensitivity}

\begin{figure}[H]
  \centering
  \includegraphics[width=\textwidth]{rep_overlay_vary_p}
  \caption{Effect of $p\in\{3/5,2/3,3/4,9/10\}$ ($M=4$, $T=10$).
           Rows: strategy frequencies.  Columns: $p$ values.
           Solid = Pop~1, dashed = Pop~2.}
  \label{fig:rep_vary_p}
\end{figure}

\begin{figure}[H]
  \centering
  \includegraphics[width=\textwidth]{rep_overlay_vary_M}
  \caption{Effect of $M\in\{2,4,6,8\}$ ($p=3/4$, $T=10$).}
  \label{fig:rep_vary_M}
\end{figure}

\begin{figure}[H]
  \centering
  \includegraphics[width=\textwidth]{rep_overlay_vary_T}
  \caption{Effect of $T\in\{2,5,10,20\}$ ($p=3/4$, $M=4$).}
  \label{fig:rep_vary_T}
\end{figure}

\subsection{Initial Condition Sensitivity}

\begin{figure}[H]
  \centering
  \includegraphics[width=\textwidth]{rep_overlay_vary_x0_p075}
  \caption{Effect of varying initial condition ($p=3/4$, $M=4$, $T=10$).
           Symmetric start: $x_0=y_0$.}
  \label{fig:rep_x0_p075}
\end{figure}

\begin{figure}[H]
  \centering
  \includegraphics[width=\textwidth]{rep_overlay_vary_x0_p060}
  \caption{Effect of varying initial condition ($p=3/5$, $M=4$, $T=10$).}
  \label{fig:rep_x0_p060}
\end{figure}

%============================================================
\section{Evolutionary Stable Strategies}
\label{sec:ESS}
%============================================================

\begin{definition}[Two-Population ESS]\label{def:ESS}
A strategy $\hat x\in\Delta_3$ for Population~1 is \emph{evolutionarily stable}
against Population~2 playing $\hat y\in\Delta_3$ iff for all
$x\neq\hat x$:
\[
  \hat x^\top C\hat y \geq x^\top C\hat y,\qquad
  \hat x^\top C\hat y = x^\top C\hat y \;\Rightarrow\;
  \hat x^\top Cy > x^\top Cy.
\]
A \emph{two-population ESS pair} $(\hat x,\hat y)$ requires this condition for
both populations simultaneously.
\end{definition}

\begin{proposition}[ESS Pairs]\label{prop:ESS}
For $M\geq 4$ and large $T$:
\begin{enumerate}
  \item The pair $((0,1,0),(0,1,0))$ is a two-population ESS iff $p>2/3$.
  \item No pure pair $(e_i,e_j)$ with $i,j\in\{1,3\}$ is a two-population ESS:
        the zero eigenvalues in Proposition~\ref{prop:stabAsym} prevent strict
        satisfaction of the ESS second condition.
  \item No interior point of $\mathcal{F}$ is an ESS.
\end{enumerate}
\end{proposition}

\begin{proof}
\textit{Part~1}: $\hat x^\top C\hat y=T$.  For $x\neq e_2$:
$x^\top Ce_2=\sum_m x_m C(m,2)$.
$C(1,2)=3(1-p)T<T$ and $C(3,2)=2-3p+T<T$ both hold iff $p>2/3$.
Hence the condition holds strictly, confirming ESS.

\textit{Part~2}: At $(e_1,e_3)$, $\hat x^\top C\hat y=C(1,3)=TM$ and
$e_3^\top Ce_3=TM=\hat x^\top Ce_3$.  The second ESS condition requires
$\hat x^\top Ce_3 > e_3^\top Ce_3$, i.e.\ $TM>TM$: false.  Hence not ESS.

\textit{Part~3}: Points in the interior of $\mathcal{F}$ are neutral to
perturbations within $\mathcal{F}$ (all earn $TM$), failing the strict ESS
condition.
\end{proof}

%============================================================
\section{Markovian Analysis}
\label{sec:markov}
%============================================================

\subsection{Two-Population Markov Chain}

The population state is $(\mathbf{s},\mathbf{t})$ where
$\mathbf{s}=(s_1,s_2,s_3)\in\mathbb{Z}_{\geq0}^3$ with $\sum_m s_m=N_1$ and
$\mathbf{t}=(t_1,t_2,t_3)\in\mathbb{Z}_{\geq0}^3$ with $\sum_n t_n=N_2$.
The state space has $\binom{N_1+2}{2}\cdot\binom{N_2+2}{2}$ elements.

Since there are no within-population games, per-player payoffs are:
\begin{align}
  q_m^{(1)}(\mathbf{s},\mathbf{t}) &= \sum_n t_n\,C(m,n),\label{eq:q1}\\
  q_n^{(2)}(\mathbf{s},\mathbf{t}) &= \sum_m s_m\,C(m,n).\label{eq:q2}
\end{align}
Note that~\eqref{eq:q1} is the payoff to a Population-1 strategy-$m$ player
against all $N_2$ Population-2 players, and~\eqref{eq:q2} is the payoff to a
Population-2 strategy-$n$ player against all $N_1$ Population-1 players.

\begin{definition}[Two-Population PPI]\label{def:PPI}
At each step, a population is chosen uniformly at random.
Within the chosen population, one player is selected uniformly and applies the
PPI revision protocol:
\[
  \varrho_{m_1 m_2}^{(1)}(\mathbf{s},\mathbf{t})
  = \frac{x_{m_2}^{(1)}\,[q_{m_2}^{(1)}-q_{m_1}^{(1)}]_+}
         {\sum_m x_m^{(1)}\,[q_m^{(1)}-q_{m_1}^{(1)}]_+},
\]
where $x_m^{(1)}=s_m/N_1$.  An analogous formula applies within Population~2.
The state transition $(\mathbf{s},\mathbf{t})\to(\mathbf{s}-\mathbf{e}_{m_1}+\mathbf{e}_{m_2},\mathbf{t})$
has probability $\tfrac{1}{2}\cdot(s_{m_1}/N_1)\cdot\varrho_{m_1m_2}^{(1)}$.
\end{definition}

\subsection{Absorbing States}

\begin{proposition}[Absorbing States]\label{prop:abs}
For $M\geq 4$, $N_1,N_2\geq 2$, and large $T$, the absorbing states of the
two-population Markov chain are:
\begin{enumerate}
  \item All states $(\mathbf{s},\mathbf{t})$ with $s_2=t_2=0$:
        \[
          \mathcal{A} = \{(\mathbf{s},\mathbf{t}):
            s_1+s_3=N_1,\; t_1+t_3=N_2,\; s_2=t_2=0\}.
        \]
        This face has $\binom{N_1+1}{1}\cdot\binom{N_2+1}{1}=(N_1+1)(N_2+1)$ states.
  \item The pure All-$1$ state $((0,N_1,0),(0,N_2,0))$.
\end{enumerate}
\end{proposition}

\begin{proof}
\textit{$\mathcal{A}$ is absorbing.}  For $(\mathbf{s},\mathbf{t})\in\mathcal{A}$
($s_2=t_2=0$): all Population-1 strategies earn $q_m^{(1)}=\sum_n t_n C(m,n)=TM$
(since $t_2=0$ and $C(m,j)=TM$ for $j\in\{1,3\}$ and $m\in\{1,3\}$).  So
all Population-1 players have equal payoff and PPI gives no switches.
Similarly for Population~2.  Hence $\mathcal{A}$ is absorbing.

\textit{$((0,N_1,0),(0,N_2,0))$ is absorbing.}  With all players using All-$1$,
there is no other strategy to imitate.

\textit{All other states are transient.}  Any state with $s_2>0$ and some
$t_m>0$ for $m\in\{1,3\}$ has Population-2 strategies earning
$q_m^{(2)}=s_2 C(2,m)+(N_1-s_2)TM$, while Population-2's All-$1$ earns
$q_2^{(2)}=s_2 C(2,2)+(N_1-s_2)C(1,2)$.
For large $T$: $q_m^{(2)}\approx(N_1-s_2)TM>q_2^{(2)}\approx N_1 T$.
So All-$1$ in Population~2 loses fitness to cooperators in Population~2, and
with positive probability Population~2's All-$1$ count decreases.  By
induction $s_2$ eventually reaches 0.
\end{proof}

\begin{remark}[Contrast with one-population case]
In the one-population ISC with All-$M$, the $s_2=0$ face is absorbing because
All-$M$ and Grim earn identical payoffs ($TM$) against each other.
In the two-population model this remains true \emph{within} each population's
internal comparisons (since $q_m^{(1)}=TM$ for all $m\in\{1,3\}$ when
$t_2=0$), so the absorbing face $\mathcal{A}$ is preserved in its full
$(N_1+1)(N_2+1)$-state extent.  This is a \emph{larger} absorbing set than the
one-population case: the two populations independently choose their level of
cooperation along the All-$M$--Grim spectrum.
\end{remark}

\subsection{State Transition Graphs (Figures~4--5)}

\begin{figure}[H]
  \centering
  \begin{subfigure}[b]{0.48\textwidth}
    \includegraphics[width=\textwidth]{mark_stategraph_p075}
    \caption{$p=3/4$.}
    \label{fig:sg_p075}
  \end{subfigure}
  \hfill
  \begin{subfigure}[b]{0.48\textwidth}
    \includegraphics[width=\textwidth]{mark_stategraph_p060}
    \caption{$p=3/5$.}
    \label{fig:sg_p060}
  \end{subfigure}
  \caption{Two-population Markov state transition graphs, $M=4$, $T=10$, $N_1=N_2=5$.
           \emph{Left of each subfigure}: Population-1 simplex projection.
           \emph{Right}: Population-2 simplex projection.
           Blue dots: states on absorbing face~$\mathcal{A}$ or pure All-$1$.
           Red dots: transient states.
           Arrow width $\propto$ transition probability.}
  \label{fig:sg_both}
\end{figure}

\subsection{Time Series}

\begin{figure}[H]
  \centering
  \begin{subfigure}[b]{0.48\textwidth}
    \includegraphics[width=\textwidth]{mark_timeseries_p075}
    \caption{$p=3/4$.}
  \end{subfigure}
  \hfill
  \begin{subfigure}[b]{0.48\textwidth}
    \includegraphics[width=\textwidth]{mark_timeseries_p060}
    \caption{$p=3/5$.}
  \end{subfigure}
  \caption{Two-population Markov dynamics time series, $M=4$, $T=10$,
           $N_1=N_2=10$, $5000$ steps, $\mathbf{s}_0=(3,4,3)$,
           $\mathbf{t}_0=(4,3,3)$.  Solid: Pop~1.  Dashed: Pop~2.}
  \label{fig:mark_timeseries}
\end{figure}

\subsection{Parameter Sensitivity}

\begin{figure}[H]
  \centering
  \includegraphics[width=\textwidth]{mark_overlay_vary_p}
  \caption{Effect of $p\in\{3/5,2/3,3/4,9/10\}$ ($M=4$, $T=10$, $N=10$).}
  \label{fig:mark_vary_p}
\end{figure}

\begin{figure}[H]
  \centering
  \includegraphics[width=\textwidth]{mark_overlay_vary_M}
  \caption{Effect of $M\in\{2,4,6,8\}$ ($p=3/4$, $T=10$, $N=10$).}
  \label{fig:mark_vary_M}
\end{figure}

\begin{figure}[H]
  \centering
  \includegraphics[width=\textwidth]{mark_overlay_vary_T}
  \caption{Effect of $T\in\{2,5,10,20\}$ ($p=3/4$, $M=4$, $N=10$).}
  \label{fig:mark_vary_T}
\end{figure}

\begin{figure}[H]
  \centering
  \includegraphics[width=\textwidth]{mark_overlay_vary_N}
  \caption{Effect of $N\in\{5,10,15,20\}$ ($N_1=N_2=N$, $p=3/4$, $M=4$, $T=10$).}
  \label{fig:mark_vary_N}
\end{figure}

\subsection{Initial State Sensitivity}

\begin{figure}[H]
  \centering
  \includegraphics[width=\textwidth]{mark_overlay_vary_s0_p075}
  \caption{Initial state sensitivity ($p=3/4$, $M=4$, $T=10$, $N=10$).}
  \label{fig:mark_s0_p075}
\end{figure}

\begin{figure}[H]
  \centering
  \includegraphics[width=\textwidth]{mark_overlay_vary_s0_p060}
  \caption{Initial state sensitivity ($p=3/5$, $M=4$, $T=10$, $N=10$).}
  \label{fig:mark_s0_p060}
\end{figure}

%============================================================
\section{Connection to Learning Theory}
\label{sec:learning}
%============================================================

Both evolutionary game theory and individual learning models converge to the
same formal dynamics under different interpretations~\cite{BorgersSarin1997}.

In the two-population context:
\begin{itemize}
  \item \textbf{Evolutionary interpretation}: Population~1 evolves against
        Population~2.  The cooperative face $\mathcal{F}$ is stable because
        any All-$1$ mutant in either population faces only cooperative opponents
        from the other population, earning $\approx T\ll TM$.
  \item \textbf{Learning interpretation}: Individual agents in Population~1
        update their strategies based on payoff comparisons with other
        Population-1 agents; their payoffs come from Population-2 opponents.
        This is the natural model for \emph{between-group learning}: players
        learn from their own group but earn from interactions with the other.
  \item \textbf{PPI and reinforcement}: Under PPI~\cite{Sandholm2010}, the
        two-population Markov chain converges in distribution (as $N_1,N_2\to\infty$)
        to the replicator dynamics~\eqref{eq:rep1}--\eqref{eq:rep2}, by an
        analogous argument to Proposition~7.4 of~\cite{LK2026}.
\end{itemize}

%============================================================
\section{Discussion}
\label{sec:disc}
%============================================================

\subsection{New Phenomena vs.\ One-Population Model}

Table~\ref{tab:compare} summarises the key differences between the
one-population and two-population ISC models.

\begin{table}[H]
\centering
\caption{Comparison: one-population vs.\ two-population ISC.}
\label{tab:compare}
\small
\renewcommand{\arraystretch}{1.2}
\begin{tabular}{lll}
\toprule
Feature & One-population & Two-population \\
\midrule
Replicator dynamics   & $\dot x_m=x_m[(Cx)_m-x^\top Cx]$ &
                        Eqs.~\eqref{eq:rep1}--\eqref{eq:rep2} \\
Accessible NE         & Symmetric only & Symmetric and asymmetric \\
Absorbing face        & 1-D ($s_2=0$ line) & 2-D ($s_2=t_2=0$ face) \\
Total absorbing states & $N+1$ & $(N_1+1)(N_2+1)$ \\
Payoff in $\mathcal{A}$ & All-$M$ earns $TM$, Grim $TM$ & Same \\
Stability of $\mathcal{F}$ & Neutral (eigenvalue = 0) & Asymptotically stable \\
$p^*$ for All-$1$ stability & $2/3$ & $2/3$ (preserved) \\
\bottomrule
\end{tabular}
\end{table}

\subsection{Asymmetric Equilibria and Their Interpretation}

The most striking new feature is the reachability of asymmetric equilibria such
as $(e_1,e_3)$ (Pop~1 all-cooperating, Pop~2 all-Grim) from generic initial
conditions.  Under~\eqref{eq:rep1}--\eqref{eq:rep2}, if Pop~2 evolves faster
toward Grim, Pop~1's All-$M$ and Grim both earn $TM$ against Pop~2's Grim,
so Pop~1 has no selection pressure to switch.  The system can lock into the
asymmetric state.

\subsection{The $p=2/3$ Threshold}

The critical threshold $p^*=2/3$ is entirely preserved:
\begin{itemize}
  \item $p>2/3$: the All-$1$ equilibrium is a two-population ESS.
  \item $p<2/3$: All-$1$ is unstable in both populations simultaneously;
        the system converges to the cooperative face $\mathcal{F}$.
\end{itemize}

%============================================================
\section{Conclusion}
\label{sec:conc}
%============================================================

We have studied the two-population ISC with strategies All-$M$, All-$1$, and
Grim.  The main results are:

\begin{enumerate}
  \item \textbf{Payoff matrix}: $C$ is the same $3\times3$ matrix as in the
        one-population model; the bimatrix game is $(C,C^\top)$.
  \item \textbf{Nash equilibria}: A continuum of asymmetric NE of the form
        $((1-a',0,a'),(1-a'',0,a''))$ exists for all $a',a''\in[0,1]$, in
        addition to all symmetric NE of the original paper.
  \item \textbf{Two-population replicator dynamics}: Equations~\eqref{eq:rep1}--\eqref{eq:rep2}
        preserve the $p^*=2/3$ stability threshold for the All-$1$ equilibrium.
        The cooperative face $\mathcal{F}$ is asymptotically stable in the
        transverse directions (stronger than the one-population result).
  \item \textbf{ESS}: All-$1$ is the unique two-population ESS, iff $p>2/3$.
  \item \textbf{Markov chain}: The absorbing face $\mathcal{A}$ is
        two-dimensional, with $(N_1+1)(N_2+1)$ states, reflecting the
        independent cooperation levels of the two populations.
  \item \textbf{Parameter sensitivity}: Results are robust to variations in
        $p$, $M$, $T$, $N$, and initial conditions.
\end{enumerate}

%------------------------------------------------------------
\begin{thebibliography}{9}

\bibitem{Rosenthal1981}
  R.\,W.~Rosenthal (1981).
  ``Games of perfect information, predatory pricing and the chain-store
  paradox.''
  \textit{Journal of Economic Theory}, 25, 92--100.

\bibitem{LK2026}
  I.~Lazaridis and A.~Kehagias (2026).
  \textit{Iterated Symmetric Centipede: Learning and Evolutionary Games}.
  Preprint, March 2026.

\bibitem{Axelrod1981}
  R.~Axelrod and W.\,D.~Hamilton (1981).
  ``The evolution of cooperation.''
  \textit{Science}, 211, 1390--1396.

\bibitem{BorgersSarin1997}
  T.~B\"{o}rgers and R.~Sarin (1997).
  ``Learning through reinforcement and replicator dynamics.''
  \textit{Journal of Economic Theory}, 77, 1--14.

\bibitem{Hofbauer1998}
  J.~Hofbauer and K.~Sigmund (1998).
  \textit{Evolutionary Games and Population Dynamics}.
  Cambridge University Press.

\bibitem{Sandholm2010}
  W.\,H.~Sandholm (2010).
  \textit{Population Games and Evolutionary Dynamics}.
  MIT Press.

\bibitem{Kemeny1969}
  J.\,G.~Kemeny and J.\,L.~Snell (1969).
  \textit{Finite Markov Chains}.
  Van Nostrand.

\bibitem{Hopkins2002}
  E.~Hopkins (2002).
  ``Two competing models of how people learn in games.''
  \textit{Econometrica}, 70, 2141--2166.

\bibitem{Cressman2003}
  R.~Cressman (2003).
  \textit{Evolutionary Dynamics and Extensive Form Games}.
  MIT Press.

\end{thebibliography}

%============================================================
\appendix
%============================================================

\section{Two-Population Replicator Dynamics: Explicit Equations}
\label{app:explicit}

Writing out~\eqref{eq:rep1}--\eqref{eq:rep2} explicitly with
$\bar\pi = x^\top Cy$:

\begin{align*}
\dot x_1 &= x_1\bigl[TM\,y_1 + 3(1-p)T\,y_2 + TM\,y_3 - \bar\pi\bigr],\\
\dot x_2 &= x_2\bigl[3pT\,y_1 + T\,y_2 + (3p+T-1)\,y_3 - \bar\pi\bigr],\\
\dot x_3 &= x_3\bigl[TM\,y_1 + (3(1-p)+T-1)\,y_2 + TM\,y_3 - \bar\pi\bigr],
\end{align*}
\begin{align*}
\dot y_1 &= y_1\bigl[TM\,x_1 + 3pT\,x_2 + TM\,x_3 - \bar\pi\bigr],\\
\dot y_2 &= y_2\bigl[3(1-p)T\,x_1 + T\,x_2 + (3(1-p)+T-1)\,x_3 - \bar\pi\bigr],\\
\dot y_3 &= y_3\bigl[TM\,x_1 + (3p+T-1)\,x_2 + TM\,x_3 - \bar\pi\bigr].
\end{align*}

\noindent These are implemented in \texttt{repDyn2pop.m}.

\section{Large-$T$ Approximation of $C$}

\[
\frac{C}{T}\approx\begin{pmatrix} M & 3(1-p) & M \\ 3p & 1 & 1 \\ M & 1 & M\end{pmatrix}.
\]

At large $T$ the off-diagonal entries $C(2,3)/T\approx 1$ and $C(3,2)/T\approx 1$
are both approximately $1$, so All-$1$ earns the same against Grim as against
itself.  The cooperative strategies both earn $M$ against each other, $M\gg 1$
for $M\geq 4$, explaining why the cooperative face is strongly attracting.

\section{MATLAB Implementation Notes}

All figures are generated by two main scripts:
\begin{itemize}
  \item \texttt{xRep2pop001.m}: integrates~\eqref{eq:rep1}--\eqref{eq:rep2}
        via \texttt{ode45} (implemented in \texttt{repDyn2pop.m}) and saves
        all replicator dynamics figures.
  \item \texttt{xMark2pop001.m}: runs the two-population Markov simulations
        via \texttt{sMarkDyn2pop.m} and builds the exact transition matrix
        via \texttt{buildMarkovP2pop.m}.
\end{itemize}

Supporting files:
\begin{itemize}
  \item \texttt{computeC.m}: computes the $3\times3$ ISC payoff matrix.
  \item \texttt{repDyn2pop.m}: ODE RHS for equations~\eqref{eq:rep1}--\eqref{eq:rep2}.
  \item \texttt{xRepDyn2pop.m}: wrapper calling \texttt{ode45}.
  \item \texttt{PhasePlot2pop.m}: two-simplex phase portrait.
  \item \texttt{ternary2cart.m}, \texttt{drawSimplex.m}: ternary coordinate
        utilities.
  \item \texttt{sMarkDyn2pop.m}: Monte Carlo Markov simulation.
  \item \texttt{buildMarkovP2pop.m}: exact transition matrix.
  \item \texttt{PStateTransitionGraph2pop.m}: state transition graph.
\end{itemize}

\end{document}
